\chapter{Ковёр: от номера метки до координат}

Ковёр --- это сетка из метках \texttt{tag36h11}, и каждая метка несёт своё число. Автопилот превращает его в положение в комнате по очень простому правилу.

\section{Номер метки равен номеру клетки}

Клеток 12 в ряду. Номер метки совпадает с номером клетки (0--143), и по нему считаются строка и столбец:

\[
\text{строка} = \lfloor id / 12 \rfloor, \qquad \text{столбец} = id \bmod 12.
\]

Шаг сетки \textbf{40 см}. Это подтверждено в бинарнике автопилота (функции деления на 12) и совпадает с нумерацией клеток в приложении. \pr{}

\begin{figure}[tbp]
\centering
\begin{tikzpicture}[x=11mm, y=11mm,
  cell/.style={draw=rule, fill=paper, minimum width=11mm, minimum height=11mm, font=\sffamily\small, inner sep=0pt},
  hl/.style={cell, fill=accentsoft, draw=accent, line width=1pt}
]
\foreach \r in {0,1,2,3}
  \foreach \c in {0,1,2,3,4,5} {
    \pgfmathtruncatemacro{\n}{\c+12*\r}
    \ifnum\n=0 \node[hl] at (\c,-\r) {\n}; \else \node[cell] at (\c,-\r) {\n}; \fi
  }
\node[font=\sffamily\small, color=inkgray] at (6.4,0) {$\cdots$};
\draw[-{Stealth}, color=accentdark, line width=0.8pt] (-0.9,0.45) -- node[above, font=\sffamily\scriptsize, color=inkgray]{столбец $\to$} (5.4,0.45);
\draw[-{Stealth}, color=accentdark, line width=0.8pt] (-0.9,0.45) -- node[left, font=\sffamily\scriptsize, color=inkgray, rotate=90, yshift=2mm]{строка $\to$} (-0.9,-3.4);
\node[font=\sffamily\scriptsize, color=inkgray, align=left] at (3,-4.1) {шаг клетки 40 см, номер = столбец + 12 $\times$ строка};
\end{tikzpicture}
\caption{Фрагмент ковра. Метка 0 --- левый верхний угол; именно её мы клали под камеру.}\label{fig:mat}
\end{figure}

\section{Положение камеры}

Кроме номера, модуль присылает смещение камеры от центра метки. Автопилот складывает положение клетки и это смещение:

\begin{tabularx}{\textwidth}{@{}l Y@{}}
\toprule
\hdr{Ось} & \hdr{Формула (метры)} \\ \midrule
Первая (по строке) & $\bigl(y \cdot 0{,}01 \cdot 8{,}333 + \text{строка} \cdot 40\bigr) \cdot 0{,}01$ \\
Вторая (по столбцу) & $\bigl(x \cdot 0{,}01 \cdot 8{,}333 + \text{столбец} \cdot 40\bigr) \cdot 0{,}01$ \\
Высота & расстояние $\cdot\, 10{,}2 \cdot$ (малый коэффициент, около $10^{-4}$) \\ \bottomrule
\end{tabularx}

\vspace{4pt}
\pr{} Константы 40, 8,333 и 10,2 лежат в сегменте данных автопилота по адресу \texttt{0x51200}. Угол приводится к диапазону $\pm 180$\textdegree.

\begin{zametka}[Проверка на живом кадре]
Для метки 0 с $x = -4$, $y = 22$ и расстоянием 296 формула даёт около $+0{,}018$ м и $-0{,}003$ м по осям и $0{,}30$ м по высоте. Дальномер показывал 287 мм. Значит, единица $x$ и $y$ --- примерно 0,83 мм, единица расстояния --- примерно 1,02 мм. \vyv{} Положение в памяти дрона вживую мы не снимали, расчёт по формуле.
\end{zametka}

\section{Что говорит образ 1.2.8}

В образе 1.2.8 формула шире: есть режим «малых» меток (для номеров больше 99, шаг 5 и масштаб 2,083), особые константы для номеров 17--20 и проверка расстояния. В старой прошивке дрона, как мы видели, такого кода нет. \pr{} Почему номера 17--20 особые, мы не знаем. \nei{}

\section{Можно ли заменить AprilTag на ArUco}

Прямой ответ: \textbf{автопилоту всё равно}. Он принимает только число и смещения, а какое семейство меток эти числа породило, ему безразлично. Но между метками и автопилотом стоит модуль.

\begin{tabularx}{\textwidth}{@{}p{4.6cm} Y@{}}
\toprule
\hdr{Что хотите} & \hdr{Как} \\ \midrule
Другое семейство AprilTag (например \texttt{tag16h5}) & Поменять одну константу в программе модуля (маска \texttt{0x10}). Таблицы остальных четырёх семейств уже в библиотеке. \vyv{} \\
ArUco & Готового кода нет. В модуле есть C++-библиотека \texttt{libopencv\_aruco}, а Python-привязки \texttt{cv2} нет: нужна программа на C++ (кросс-компилятор под ARM). Её задача --- слать кадр из главы~\ref{ch:uart}. \vyv{} \\
Другая сетка (не 12 колонок) & Константа 12 зашита в автопилоте; её пришлось бы менять там. \vyv{} \\ \bottomrule
\end{tabularx}

\vspace{4pt}
\begin{vnimanie}[Размер меток]
Размер метки и калибровка камеры влияют на расстояние и смещение. Где в программе модуля лежит калибровка, мы пока не выяснили, поэтому «заменить метки» нельзя считать готовым рецептом. \nei{}
\end{vnimanie}
